%% Algorithmic Redistricting Portfolio Summary
%% 6-page research statement with detailed paper descriptions

\documentclass[11pt]{article}

% Packages
\usepackage[utf8]{inputenc}
\usepackage[margin=1in]{geometry}
\usepackage{times}
\usepackage{amsmath}
\usepackage{booktabs}
\usepackage{hyperref}
\usepackage{enumitem}

% Title
\title{\textbf{Algorithmic Redistricting Portfolio}\\
\large An 11-Paper Research Program}

\author{Giovanni Della-Libera}
\date{February 2026}

\begin{document}

\maketitle

\section*{Overview}

Congressional redistricting in the United States suffers a legitimacy crisis: partisan mapmakers manipulate district boundaries to predetermine election outcomes, eroding public confidence in representative democracy. This 11-paper research program demonstrates that \textbf{recursive graph bisection}---a computational method used for supercomputer workload distribution---provides an algorithmic alternative that \emph{cannot gerrymander because it cannot access partisan data}.

Applied to all 50 states across three census decades (2000/2010/2020), this method produces 435 congressional districts with near-perfect population equality, geographic contiguity, and geometric compactness. Algorithmic plans create \textbf{69 more majority-minority districts} than enacted maps, exceeding Voting Rights Act requirements without explicit racial targeting. We identify a \textbf{42\% demographic threshold} determining where proportional minority representation becomes geographically feasible. Efficiency gap analysis reveals \textbf{62\% partisan bias reduction} compared to enacted plans, establishing empirical benchmarks for neutral redistricting.

\subsection*{Four Core Contributions}

\begin{itemize}[leftmargin=*,noitemsep,topsep=4pt]
    \item \textbf{Philosophical}: \emph{Impossibility defense}---algorithm cannot see partisan data, therefore cannot gerrymander
    \item \textbf{Technical}: Edge-weighted graph partitioning achieves +56\% compactness improvement while enabling VRA compliance
    \item \textbf{Empirical}: National-scale validation across 50 states and 3 census decades (150 redistricting scenarios)
    \item \textbf{Legal}: 42\% state minority threshold determines VRA feasibility; +69 majority-minority district surplus; 62\% partisan bias reduction
\end{itemize}

\subsection*{Portfolio Architecture}

The 11-paper portfolio follows a hierarchical structure reflecting methodological dependencies:

\begin{itemize}[leftmargin=*,noitemsep,topsep=4pt]
    \item \textbf{Foundation} (Paper 01): Recursive bisection method and impossibility defense
    \item \textbf{Primary Extensions} (Papers 02-03): Compactness optimization and demographic awareness
    \item \textbf{Systematic Investigations} (Papers 04-10): Seven papers exploring properties, limits, and alternatives
    \item \textbf{Partisan Fairness} (Paper 11): Efficiency gap analysis quantifying bias reduction
\end{itemize}

This structure enables progressive deepening: Papers 04-11 build on the compactness optimization (02) and VRA framework (03), enabling systematic investigation of algorithm properties and political implications.

\section*{Paper Summaries}

All 11 papers are complete and ready for submission.

\subsection*{Paper 01: Recursive Bisection for Congressional Redistricting}

This foundation paper establishes the core method and philosophical framework. We demonstrate that recursive graph bisection---the same algorithm used to distribute workloads across supercomputer processors---can partition census tracts into congressional districts meeting all constitutional requirements. The method operates exclusively on population counts and geographic adjacency, structurally unable to access partisan voter data. This creates an \emph{impossibility defense}: the algorithm cannot gerrymander not because it chooses not to, but because it lacks the information required to do so. Applied to all 50 states, the method produces 435 districts with population equality within 0.5\%, full geographic contiguity, and Polsby-Popper compactness scores 23\% higher than enacted plans. The paper frames algorithmic redistricting as extending the Huntington-Hill precedent for congressional apportionment (successful since 1941) to the boundary design problem, demonstrating that mathematical governance can restore electoral legitimacy through structural objectivity rather than institutional reform.

\subsection*{Paper 02: Edge-Weighted Graph Partitioning for Compact Districts}

This technical paper introduces edge-weighting to optimize district compactness. Standard graph bisection minimizes the number of edges cut between partitions but treats all edges equally, producing districts that may have irregular boundaries. We assign higher costs to edges between geographically similar census tracts (based on shared boundary length and demographic similarity), making these edges expensive to cut and encouraging compact, spatially cohesive districts. Applied nationally, edge-weighting with scaling parameter $\alpha=5$ improves Polsby-Popper compactness by 56\% over unweighted partitioning (0.39 vs 0.25 mean) while maintaining population equality. The method is computationally efficient (O($n$ log $k$) for $k$ districts) and requires no human judgment about ``desirable'' district shapes. This establishes that purely algorithmic methods can achieve compactness levels comparable to commission-drawn maps (mean PP = 0.42) without manual boundary adjustments.

\subsection*{Paper 03: VRA Compliance Without Explicit Racial Targeting}

This legal-empirical paper demonstrates that algorithmic redistricting exceeds Voting Rights Act requirements without using racial data as input. We modify edge-weighting to assign higher costs to edges connecting census tracts with high minority voting-age population (VAP), encouraging these tracts to remain together in the same districts. This creates majority-minority (MM) districts where minority voters $\geq$ 50\% VAP, satisfying Section 2 of the VRA. Critically, the algorithm receives only demographic counts (``Tract A has 2,400 Black VAP''), not geographic coordinates or partisan voting history. Applied to all 50 states in 2020, algorithmic plans create 137 MM districts compared to 68 in enacted plans---a surplus of +69 districts (+101\%). Key examples: Alabama (0$\rightarrow$2 MM districts), Georgia (5$\rightarrow$6), Louisiana (1$\rightarrow$2). This surplus demonstrates that VRA compliance does not require partisan manipulation or manual map-drawing; purely algorithmic methods can produce more MM districts than human mapmakers while maintaining geographic compactness and population equality.

\subsection*{Paper 04: The 42\% Threshold for VRA Feasibility}

This legal-theoretical paper identifies a demographic threshold determining where proportional minority representation becomes geographically feasible. We analyze the relationship between state minority population percentage and MM district creation across all 50 states and three census decades. States above 42\% minority VAP achieve near-proportional outcomes (e.g., Mississippi at 44\% minority creates 2 of 4 MM districts, matching 50\% proportionality). States below 37\% face geometric constraints that prevent proportional representation even with aggressive VRA-weighted partitioning (e.g., Wisconsin at 18\% minority cannot create MM districts through geographic clustering alone). This 42\% threshold has critical legal implications: it operationalizes the first \emph{Gingles} precondition (``sufficiently large and geographically compact minority population'') with quantitative precision. Courts evaluating VRA claims can use this threshold to assess whether proportional MM district creation is geographically feasible or structurally impossible in a given state. The finding also explains why some states face unavoidable tension between the VRA and compactness criteria: proportional representation requires geographic concentration that may not naturally exist.

\subsection*{Paper 05: Cross-Census Validation via Slice-Based Methodology}

This methodological paper validates algorithmic redistricting across census decades using tract correspondence analysis. Census tract boundaries change between decades, making direct district-level comparison difficult (2000 tracts $\neq$ 2010 tracts). We develop a slice-based methodology: identify tract splits/merges, create correspondence weights, and propagate district assignments across decades. This enables three-way comparison: same geographic region, three different census snapshots. Variance decomposition reveals that differences between geographic regions (e.g., urban vs. rural) exceed temporal changes by 3.2×, indicating districts reflect durable spatial structure rather than decade-specific population artifacts. The method provides evidence of algorithmic stability: recursive bisection maintains consistent spatial patterns despite 10-year demographic shifts. This addresses a key criticism of automated redistricting (``arbitrary boundaries that change drastically each decade'') by demonstrating empirically that algorithmic plans exhibit geographic continuity comparable to human-drawn maps. The slice methodology generalizes beyond redistricting to any geographic analysis requiring temporal comparison across changing boundaries.

\subsection*{Paper 06: VRA-Compactness Pareto Frontiers}

This optimization paper challenges the conventional wisdom that VRA compliance requires sacrificing compactness. We generate Pareto frontiers for representative states (Alabama, Georgia, Louisiana) by systematically varying the edge-weighting parameter $\alpha$ (compactness emphasis) and demographic clustering strength $\beta$ (VRA emphasis). Counter-intuitively, non-MM districts gain +7.5\% compactness when VRA-weighting is enabled, not lose compactness. This occurs because demographic clustering concentrates minority voters into specific districts, freeing the algorithm to optimize compactness more aggressively in remaining districts. The analysis reveals three regions of the parameter space: (1) VRA-optimal ($\beta \geq 10$): maximizes MM districts, slight compactness reduction; (2) balanced ($5 \leq \beta < 10$): near-maximum MM districts with high compactness; (3) compactness-optimal ($\beta < 5$): maximum compactness but may fail VRA requirements. Most states achieve Pareto-optimal outcomes in the balanced region, demonstrating that VRA compliance and compactness are not fundamentally in conflict when using algorithmic methods.

\subsection*{Paper 07: Temporal Stability Across Census Cycles}

This temporal analysis paper demonstrates that recursive bisection maintains higher district boundary stability across census cycles than flat n-way partitioning. We measure census tract retention: what percentage of tracts remain in the same district number from 2010 to 2020 redistricting? Recursive bisection achieves 80\% tract retention compared to 66\% for n-way partitioning, a 14 percentage point advantage. This occurs because recursive bisection creates a hierarchical scaffold (first split: North vs South; second split: NE/NW/SE/SW; etc.) that persists even as population shifts within regions. N-way partitioning solves for all $k$ districts simultaneously, producing globally optimal but temporally unstable solutions. The stability advantage has practical implications: higher tract retention means voters are more likely to remain in the same district across decades, reducing disruption to community representation and incumbent-constituent relationships. This addresses reform commission concerns that algorithmic methods produce ``chaotic'' redistricting with complete boundary redrawing every decade. Empirically, recursive bisection boundary stability matches that of human-drawn maps (78-82\% retention) while maintaining all constitutional requirements.

\subsection*{Paper 08: Why Single-Objective Edge-Weighting Outperforms Multi-Constraint}

This optimization theory paper explains why single-objective edge-weighting produces better results than multi-constraint optimization for redistricting. Multi-constraint methods (e.g., METIS \texttt{PartKway} with compactness constraint, population constraint, and VRA constraint) suffer from constraint conflicts: tightening one constraint forces violation of others. We demonstrate mathematically that the redistricting problem exhibits \emph{constraint non-commutativity}: order of constraint application affects solution quality. Edge-weighting reformulates multiple constraints as a single objective (minimize weighted edge cut), eliminating conflicts. Empirically, edge-weighting produces districts with 18\% better compactness and 12\% more MM districts than multi-constraint methods given equal population balance tolerance. The paper generalizes to other graph partitioning applications: whenever constraints interact non-linearly, reformulating as a single weighted objective often outperforms explicit multi-constraint optimization. This has implications for algorithm design beyond redistricting: load balancing, network segmentation, and VLSI design all benefit from single-objective reformulation rather than adding explicit constraints.

\subsection*{Paper 09: Parameter Sensitivity and Tree Structure Irrelevance}

This sensitivity analysis paper tests how recursive bisection outcomes depend on the hierarchical tree structure used for splitting. Does it matter whether we first split East-West vs. North-South? Does random tree initialization produce wildly different districts? We generate 100 random tree structures per state and measure district outcome variance across metrics (compactness, MM districts, population balance). With edge-weighting $\alpha \geq 5$, tree structure choice accounts for only 3\% of outcome variance---statistically insignificant. This means the algorithm is robust: different splitting orders converge to similar district configurations because edge weights guide the optimization toward consistent geographic patterns. Without edge-weighting ($\alpha=0$), tree structure accounts for 24\% of variance, indicating arbitrary outcomes. The finding addresses a key criticism of recursive methods (``outcomes depend on arbitrary tree choices'') by demonstrating empirically that edge-weighting eliminates this sensitivity. Adaptive bisection strategies (choosing split directions based on state geography) provide no additional benefit when edge-weighting is enabled, simplifying implementation: random tree initialization suffices.

\subsection*{Paper 10: N-Way vs Recursive: Statistical Equivalence}

This comparative methods paper tests whether recursive bisection and n-way partitioning (flat, non-hierarchical) produce statistically different outcomes for redistricting. We apply both methods to all 50 states with identical population constraints and edge-weighting parameters. Statistical tests reveal no significant differences in compactness (recursive: 0.39, n-way: 0.38, p=0.23), MM district counts (recursive: 137, n-way: 134, p=0.18), or population balance (both within 0.5\%). However, recursive bisection achieves 14 percentage points higher temporal stability (80\% vs 66\% tract retention) and runs 40\% faster (O($n$ log $k$) vs O($n$ log$^2$ $k$) complexity). This establishes that both methods are viable for neutral redistricting, but recursive bisection offers practical advantages for real-world implementation. The statistical equivalence finding is surprising: theory predicts recursive methods might produce suboptimal solutions due to hierarchical constraints, but edge-weighting compensates for this limitation. The paper concludes that architectural choice (recursive vs n-way) matters less than optimization design (edge-weighting vs unweighted).

\subsection*{Paper 11: Partisan Fairness via Efficiency Gap Analysis}

This partisan fairness paper quantifies bias reduction through national-scale efficiency gap analysis. The efficiency gap (EG) measures partisan asymmetry: percentage of ``wasted votes'' for one party minus wasted votes for the other party. Enacted plans show +5.1\% EG favoring Republicans (averaged across 50 states, 2016-2020 elections). Algorithmic plans show -3.2\% EG favoring Democrats, a 62\% reduction in absolute bias. Regional analysis reveals largest improvements in the Rust Belt (10 percentage point EG reduction) due to elimination of aggressive Republican gerrymandering. The persistent -3.2\% Democratic bias in algorithmic plans reflects unavoidable geographic sorting: Democrats concentrate in urban areas, producing some wasted votes even under neutral districting. This establishes an empirical lower bound: -3.2\% represents the minimum partisan asymmetry achievable under compact districting constraints. Enacted plans exceeding algorithmic baselines by $>$7 percentage points (e.g., Wisconsin at +13\% vs algorithmic -2\%, difference of 15 points) provide quantitative evidence of partisan manipulation. The paper frames algorithmic redistricting as providing neutral benchmarks for courts to evaluate gerrymandering claims post-\emph{Rucho}, establishing that substantial partisan bias reduction is achievable without explicit fairness criteria.

\section*{Portfolio-Level Themes}

\subsection*{Impossibility Defense (Papers 01, 11)}
The core philosophical innovation: algorithms cannot gerrymander because they structurally lack access to partisan data. Paper 01 establishes this principle theoretically; Paper 11 validates it empirically by demonstrating 62\% bias reduction without any partisan fairness objective in the algorithm specification.

\subsection*{VRA Compliance Without Racial Targeting (Papers 03, 04, 06)}
Three papers demonstrate that demographic awareness through edge-weighting produces more MM districts than enacted plans (+69 surplus) without explicit racial quotas. Paper 03 establishes the method, Paper 04 identifies the 42\% threshold for feasibility, and Paper 06 debunks VRA-compactness tradeoff myths.

\subsection*{Algorithmic Robustness (Papers 05, 07, 08, 09, 10)}
Five papers validate that recursive bisection produces stable, consistent, and optimal outcomes. Papers 05 and 07 establish temporal stability, Papers 08 and 09 demonstrate insensitivity to implementation choices, and Paper 10 confirms statistical equivalence with alternative methods.

\subsection*{Technical Innovation (Papers 02, 08, 09)}
Three papers advance graph partitioning methodology beyond redistricting applications. Edge-weighting (Paper 02) reformulates multi-objective optimization, constraint conflict analysis (Paper 08) applies to general partitioning problems, and sensitivity analysis (Paper 09) establishes robustness bounds.

\section*{Impact and Implications}

\subsection*{For Reform Commissions}
This portfolio provides a complete implementation blueprint: Papers 01-02 specify the algorithm, Papers 03-06 address legal compliance (VRA, compactness), Papers 07-10 validate stability and robustness, and Paper 11 quantifies partisan fairness improvements. Reform commissions can adopt algorithmic redistricting with confidence that all legal requirements are satisfied and outcomes are demonstrably neutral.

\subsection*{For Courts}
Post-\emph{Rucho}, federal courts lack judicially manageable standards for partisan gerrymandering claims. Papers 01 and 11 provide quantitative benchmarks: enacted plans deviating from algorithmic baselines by $>$7 percentage points (efficiency gap) or $>$50 MM districts suggest manipulation. This operationalizes the impossibility defense as an evidentiary standard.

\subsection*{For Scholars}
The portfolio advances three research frontiers: (1) computational political science (applying graph algorithms to electoral design), (2) optimization theory (single-objective reformulation vs multi-constraint), and (3) legal empiricism (quantifying VRA feasibility thresholds, partisan bias bounds). Each paper contributes individually to these literatures while collectively demonstrating feasibility of algorithmic governance.

\section*{Glossary of Key Terms}

\begin{description}[leftmargin=2.5cm,style=nextline,noitemsep,topsep=4pt]
    \item[Census tracts] Geographic units used as atomic partitioning units ($\sim$4,000 residents each, 74,001 tracts nationally in 2020)
    \item[METIS] Multilevel graph partitioning library developed by Karypis \& Kumar, widely used in high-performance computing
    \item[Edge-weighted partitioning] Assigning higher costs to edges between similar nodes (making them expensive to cut), encouraging spatially compact districts
    \item[Efficiency gap (EG)] Partisan asymmetry measure: (wasted votes for party A - wasted votes for party B) / total votes. Range: [-100\%, +100\%]
    \item[Majority-minority (MM) districts] Districts where minority voters $\geq$ 50\% of voting-age population (VAP), satisfying VRA Section 2
    \item[Polsby-Popper (PP)] Compactness measure: $4\pi \times \text{Area} / \text{Perimeter}^2$. Range: [0,1], higher = more compact. Circle = 1.0, typical districts = 0.3-0.5
    \item[Recursive bisection] Hierarchical method splitting graph into two parts repeatedly (binary tree structure). Complexity: O($n$ log $k$) for $k$ districts
    \item[Impossibility defense] Algorithm cannot gerrymander by design: it structurally cannot access partisan data, unlike commission-drawn maps
    \item[VRA] Voting Rights Act (1965)---federal law prohibiting practices that dilute minority voting strength. Section 2 requires proportional representation where geographically feasible
\end{description}

\section*{Code, Data, and Replication}

\textbf{Replication materials} include the \texttt{redist} Rust CLI (redist-metis crate) and sample datasets for Vermont and Delaware (smallest states, $\sim$2 minute runtime on commodity hardware). Full 50-state redistricting requires $\sim$40GB census data (U.S. Census Bureau redistricting files for 2000/2010/2020) and $\sim$18 minutes parallel for all 50 states (\texttt{redist states --workers 8}). Binary available at \url{https://github.com/giodl73-repo/REDIST}.

\textbf{Census data sources}: U.S. Census Bureau redistricting files (PL 94-171), TIGER/Line shapefiles for geographic boundaries, American Community Survey (ACS) for demographic data.

\textbf{METIS library}: \url{http://glaros.dtc.umn.edu/gkhome/metis/metis/overview}

\textbf{Contact}: giodl@microsoft.com

\section*{Citation Guide}

\textbf{Publication status}: All 11 papers are complete and ready for submission (as of February 2026). Preprints will be available on arXiv upon initial submission.

\textbf{For the entire research program}:

\begin{quote}
Della-Libera, G. (2026). Algorithmic redistricting: An 11-paper research program. Preprint series.
\end{quote}

\textbf{For specific papers} (examples):

\begin{quote}
Della-Libera, G. (2026). Recursive bisection for congressional redistricting. Preprint, submitted to \emph{American Political Science Review}.
\end{quote}

\begin{quote}
Della-Libera, G. (2026). Measuring partisan fairness in algorithmic redistricting: A national efficiency gap analysis. Preprint, submitted to \emph{American Political Science Review}.
\end{quote}

\vspace{12pt}

\noindent\rule{\textwidth}{0.4pt}

\vspace{6pt}

\noindent\textbf{About this summary}: This 6-page research statement provides detailed descriptions of the 11-paper algorithmic redistricting portfolio. For a brief overview, see the 3-page \emph{Guide to the Algorithmic Redistricting Portfolio}. For full findings, see individual papers. For synthesis across papers, see Paper 00 (Synthesis Metapaper for \emph{Science}).

\end{document}
